% !TeX root = ../navier-stokes-manuscript.tex
% ============================================================
% 2.3 The Pressure-Complete Mixed Jet
% ============================================================

\subsection{The Pressure-Complete Mixed Jet}\label{sub:ThePressureCompleteMixedJet}

Keep the temporal notation from \cref{sub:SingularityAndTheDerivativeTower}:
\[
  V_n
  :=
  \partial_t^nu,
  \qquad
  Q_n
  :=
  \partial_t^np.
\]
The temporal recurrence already derived there is
\[
  V_{n+1}
  +
  \sum_{a=0}^{n}
  \binom{n}{a}
  (V_a\cdot\nabla)V_{n-a}
  +\nabla Q_n
  =
  \nu\Delta V_n,
  \qquad
  \nabla\cdot V_n=0.
\]
At the same temporal depth, incompressibility fixes the pressure row by
\[
  -\Delta Q_n
  =
  \partial_i\partial_j
  \sum_{a=0}^{n}
  \binom{n}{a}
  (V_a)_i(V_{n-a})_j.
\]
Rapid decay fixes its normalization:
\[
  Q_n
  =
  R_iR_j
  \sum_{a=0}^{n}
  \binom{n}{a}
  (V_a)_i(V_{n-a})_j.
\]
Each temporal row still mixes time with spatial differentiation through transport, viscosity, and pressure.
To expose those derivatives as coordinates of their own, first differentiate one transport factor in a coordinate direction \(x_\ell\):
\[
  \partial_\ell\bigl((V_a\cdot\nabla)V_{n-a}\bigr)
  =
  ((\partial_\ell V_a)\cdot\nabla)V_{n-a}
  +
  (V_a\cdot\nabla)\partial_\ell V_{n-a}.
\]
At every further spatial order, the product rule divides the derivatives between the same two factors.
For a multi-index \(\alpha\in\Nzero^3\), define
\[
  J_{n,\alpha}
  :=
  \partial_t^n\partial_x^\alpha u,
  \qquad
  \Pi_{n,\alpha}
  :=
  \partial_t^n\partial_x^\alpha p.
\]
For \(\beta\leq\alpha\), write
\[
  \binom{\alpha}{\beta}
  :=
  \prod_{k=1}^{3}
  \binom{\alpha_k}{\beta_k}.
\]
The spatial product rule applied to each temporal split gives
\[
  \partial_t^n\partial_x^\alpha
  \bigl((u\cdot\nabla)u\bigr)
  =
  \sum_{a=0}^{n}
  \sum_{\beta\leq\alpha}
  \binom{n}{a}
  \binom{\alpha}{\beta}
  (J_{a,\beta}\cdot\nabla)
  J_{n-a,\alpha-\beta}.
\]
Substitution gives the pressure-complete mixed recurrence
\begin{align*}
  J_{n+1,\alpha}
  &+
  \sum_{a=0}^{n}
  \sum_{\beta\leq\alpha}
  \binom{n}{a}
  \binom{\alpha}{\beta}
  (J_{a,\beta}\cdot\nabla)
  J_{n-a,\alpha-\beta}
  \\
  &+
  \nabla\Pi_{n,\alpha}
  =
  \nu\Delta J_{n,\alpha},
  \qquad
  \nabla\cdot J_{n,\alpha}=0,
\end{align*}
with
\[
  -\Delta\Pi_{n,\alpha}
  =
  \partial_i\partial_j
  \sum_{a=0}^{n}
  \sum_{\beta\leq\alpha}
  \binom{n}{a}
  \binom{\alpha}{\beta}
  (J_{a,\beta})_i
  (J_{n-a,\alpha-\beta})_j.
\]
Every coordinate in this family is a derivative of the same \((u,p)\), and every row retains the transport, viscous, divergence, and pressure relations that produced it.
The recurrence shows how repeated time differentiation keeps opening higher spatial derivatives through the same equation.
The critical height gives a scalar calculation in which that climb can be followed one row at a time.
